% ID: SEL-D0043-OA-E01
% ESTADO: validada. S044, cotejada con examen y criterios.
\textbf{a)} Para $x\ge-1$ se exige
$\max(0,x-1)\le y\le(3-x)/2$.
Sobre $x=-1$ quedan $( -1,0)$ y $(-1,2)$.
En el eje horizontal la cota $y\ge x-1$ da el vértice $(1,0)$.
Las dos fronteras oblicuas se cruzan al resolver
$x-1=(3-x)/2$, lo que da $(5/3,2/3)$.
\begin{center}
\setlength{\unitlength}{12mm}
\begin{picture}(7,4)(-2,-1)
\qbezier(-1.6,0)(0.6,0)(2.8,0)\put(2.9,0){$x$}
\qbezier(0,0)(0,1.15)(0,2.3)\put(0,2.4){$y$}
\qbezier(-1,0)(-1,1)(-1,2)
\qbezier(-1,2)(0.33333,1.33333)(1.66667,0.66667)
\qbezier(1.66667,0.66667)(1.33333,0.33333)(1,0)
\put(-1,2.15){\makebox(0,0)[b]{$(-1,2)$}}
\put(1.8,0.7){$(5/3,2/3)$}
\put(-1,-0.15){\makebox(0,0)[t]{$(-1,0)$}}
\put(1,-0.15){\makebox(0,0)[t]{$(1,0)$}}
\end{picture}
\end{center}
\textbf{b)} $F=2(x+2y)\le6$. El máximo $\boxed{6}$
se alcanza en el lado $x+2y=3$, con $-1\le x\le5/3$.
Por otro lado, $x\ge-1$ e $y\ge0$ implican $F\ge-2$.
La igualdad exige ambas cotas a la vez, así que el mínimo
$\boxed{-2}$ se alcanza sólo en $(-1,0)$.
